\section{Results and Discussion}

\Cref{CER} reports the overall CER, we get the CER for a five-vowel inventory,
a relaxed (\emph{e}, \emph{i}, \emph{y}) case, a relaxed (\emph{o}, \emph{u}) case,
and a combination of both relaxed cases resulting to a three-vowel inventory.

\begin{table}[h]
    \centering
    \caption{Overall CER (\%) under four vowel-equivalence conditions.}
    \label{CER}
    \begin{tabular}{lcccc}
        \hline
        \textbf{System} & \textbf{5-vowel} & \textbf{\emph{e} = \emph{i} = \emph{y}} & \textbf{\emph{o} = \emph{u}} & \textbf{3-vowel} \\
        \hline
        Direct      & ---              & ---                                     & ---                          & ---              \\
        IPA-guided  & 4.67             & 3.70                                    & 4.58                         & 3.61             \\
        \hline
    \end{tabular}
\end{table}

Under the strict five-vowel condition, the system achieved a CER of 4.67\%.
Relaxing the (\emph{e}, \emph{i}, \emph{y}) equivalence reduced the CER to
3.70\%, while the (\emph{o}, \emph{u}) relaxation yielded a fairly marginal
reduction to 4.58\%. The close range of error rates suggest that remaining
character errors are not majorly attributed to the (\emph{e}, \emph{i},
\emph{y}) or (\emph{o}, \emph{u}) equivalences, contributing minimally to
overall error. Combining both equivalences produced the lowest CER at 3.61\%,
producing a 1.06 percentage-point drop.

Because English is a stress-timed language, some pronunciations get reduced to
the schwa /\textipa{@}/ sound. This creates a bit of an ambiguity between vowel
sounds as it is a very neutral sound. This results in \emph{e}, \emph{y}, and
\emph{i} errors being the most frequent ones. This may explain why relaxing the
    {\emph{e}, \emph{i}, \emph{y}} equivalence significantly reduces the CER.

The phoneme correspondence for vowels is most of the time ambiguous, splitting
a vowel pronunciation into either its long or short version. Some examples:

\begin{table}[h]
    \centering
    \caption{Examples of long and short vowel correspondence.}
    \label{vowel_examples}
    \begin{tabular}{lll}
        \hline
                 & \textbf{Long vowel} & \textbf{Short vowel} \\
        \hline
        \emph{a} & ``cake'' (``keyk'') & ``cat'' (``kat'')    \\
        \emph{e} & ``me'' (``mi'')     & ``bed'' (``bed'')    \\
        \emph{i} & ``time'' (``taym'') & ``sit'' (``sit'')    \\
        \emph{o} & ``home'' (``hom'')  & ``hot'' (``hat'')    \\
        \emph{u} & ``cube'' (``kyub'') & ``cup'' (``kap'')    \\
        \hline
    \end{tabular}
\end{table}

\begin{table}[h]
    \centering
    \caption{Comparison of Actual and Expected Nativized Forms of English Loanwords}
    \label{tab:actual_expected}
    \begin{tabular}{llll}
        \hline
        \textbf{English Word} & \textbf{IPA} & \textbf{Actual} & \textbf{Expected} \\
        \hline
        ``abortion''          & \textipa{/@"bOrS@n/}       & ``aborsyon''    & ``aborsyon''      \\
        ``escalator''         & \textipa{/"Esk@leIt@r/}    & ``eskaleytor''  & ``eskaleytor''    \\
        ``gimmick''           & \textipa{/"gImIk/}         & ``gimik''       & ``gimik''         \\
        ``mouse''             & \textipa{/maUs/}           & ``maws''        & ``maws''          \\
        ``candy''             & \textipa{/"k\ae ndi/}      & ``k\textbf{a}ndi'' & ``k\textbf{e}ndi'' \\
        ``mall''              & \textipa{/mOl/}            & ``m\textbf{a}l''   & ``m\textbf{o}l''   \\
        ``bus''               & \textipa{/bVs/}            & ``b\textbf{a}s''   & ``b\textbf{u}s''   \\
        ``football''          & \textipa{/"fUtbOl/}        & ``futb\textbf{a}l''& ``futb\textbf{o}l''\\
        \hline
    \end{tabular}
\end{table}

This ambiguity results in a problem where a vowel must be mapped to multiple
different pronunciations.
